% Pursuing Stacks (cotes 134-2 to 134-8): archive cote 134-5, site page 121 (PDF page 122), typescript page 308.
% typescript. Draft compared: ps5.tex:815-869.
% AI-transcribed by Claude Opus 5.5 (claude-opus-5-5, adaptive thinking, effort high), Message Batches API, batch ps3 (msgbatch_01XdFmuVVBAcgc7EFxdtCqsj), 2026-11-12; not yet checked by eye.
\scanpage{134-5/121}
\mspage{308}

visualize best maybe by writing
\[ M^{A} = (\underline{Hom}(A^{o},M^{o}))\over{\unsure{x}}{^{o}} \quad , \]
i.e. interpreting the dual of $M^{A}$ as the category of A-objects of $M^{o}$ ,
for instance (if $A = \hand{\Delta}$) \inserted{as} the \xout{\illegible{1}} dual of the category of ss objects of
$M^{o}$ . Now Quillen has given handy conditions, in terms of projectives
of $M^{o} = N$ namely in terms of injectives of M, for the category
$\underline{Hom}(\hand{\Delta}^{o},N)$ of ss objects of N to be a closed model category - hence
the dual category $M^{A}$ will turn out as \unsure{a} closed model category too, un-%
der suitable conditions involving existence of injective objects in M.
These conditions are satisfied for instance when M is a topos, and no-%
tably when M is of the type $B^{\wedge}$ - quite an interesting particular case
indeed ! Assuming that M is stable under both types of limits, so is
$M^{A}$ , hence there is an initial and final object, and according to
Quillen's factorization axiom, the \xout{\unsure{canonical}} map from the former to
the latter can be factored through an object
\[ F \text{ in } M^{A} \ , \quad \text{i.e.} \quad F \colon A \longrightarrow M \ , \]
which is \uline{cofibering}, and such that $F \longrightarrow e$ is a \uline{trivial fibration}.
The idea is that these conditions mean more or less, at any rate imply,
that F is a test functor.

I hit upon this ``way out'', while trying to construct test func-%
tors from $\hand{\Delta}$ to any elementary modelizer $B^{\wedge}$ , in order to try and check
that $B^{\wedge}$ is a (semi)simplicial model category in the sense of Quillen.
The intuitive idea of constructing inductively the components $F_{n}$ of F
was simple enough, still I got stuck in some messyness and did not try
to push through this way, all the \xout{more}\inserted[typed]{less} as this naive approach had no
chance of generalizing to the case of a more or less general test cate-%
gory A. Of course, for the time being Quillen's theorems\hand{,} about
certain categories $\underline{Hom}(A^{o},N)$ being closed model categories, is equal-%
ly restricted to the case when $A = \hand{\Delta}$ , which looks as usual like a
rather arbitrary assumption. Thus, to ``test'' whether the feeling \xout{abou}
about \uline{any} test category more or less being ``just as good'' as $\hand{\Delta}$ , a
second point would be to see whether Quillen's theorems extend, which
presumably is going to be very close to the first point I raised.

I take this occasion to raise a third point - where there is no
reason to restrict to an M which is modelizing (neither was there such
reason before, when phrasing everything in terms of aspheric or c-as-%
pheric functors, rather than test functors ...). Namel\over{\unsure{x}}{y}, assuming as

\margin{121}
